\section*{Chapitre \ref{ch:rc:forward-rate-and-market-models} --- Modèles de taux forward et de marché}

\begin{solution}{exo:rc:forward-rate-and-market-models:1}
$\alpha(t,T) = \sigma\int_t^T\sigma\,du = \sigma^2(T-t)$. Il s'agit du \omterm{def:rc:short-rate-models:hw}{modèle de Hull--White} sans
retour à la moyenne ($\kappa=0$), connu sous le nom de modèle de Ho--Lee~: chaque forward bouge de
le même choc normal.
\end{solution}

\begin{solution}{exo:rc:forward-rate-and-market-models:2}
Chaque forward est une martingale sous sa propre mesure de paiement, avec une volatilité
déterministe, donc il est lognormal sous celle-ci et son caplet relève de la formule de Black. Un swap
taux de swap est une somme pondérée de plusieurs forwards à poids stochastiques~; il n'est pas
lognormal sous sa probabilité annuité, donc Black n'en est qu'une approximation.
\end{solution}

\begin{solution}{exo:rc:forward-rate-and-market-models:3}
$e^{-0.4\times8} = 0.0408$~; avec $\beta=0.02$, $e^{-0.16} = 0.852$.
\end{solution}

\begin{solution}{exo:rc:forward-rate-and-market-models:4}
À $t=0.5$, la prochaine fixation est $T_1$, donc $m(t)=1$~:
$\mu_2 = \sigma_2F_2\bigl(\frac{\delta\rho_{21}\sigma_1F_1}{1+\delta F_1}+\frac{\delta\sigma_2F_2}{1+\delta F_2}\bigr)$~;
$F_1$ et $F_2$ interviennent ($F_0$ a déjà été fixé).
\end{solution}

\begin{solution}{exo:rc:forward-rate-and-market-models:5}
$\sigma_s^2 = \frac{1}{S^2}\sum w_iw_jF_iF_j\rho_{ij}\sigma^2 = 0.25\times0.04\times(2+2\rho)$~:
20.00~\% pour $\rho=1$, 17.32~\% pour $\rho=0.5$.
\end{solution}

\begin{solution}{exo:rc:forward-rate-and-market-models:6}
Un swap d'un an qui démarre dans cinq ans n'est qu'un unique forward annuel, $F_5$~: sa
swaption est un caplet, évalué par sa propre volatilité, que les deux modèles calibrent
sur la même cible~; la corrélation n'a d'effet qu'à partir de deux forwards au moins.
\end{solution}

\begin{solution}{exo:rc:forward-rate-and-market-models:7}
2\,351\,823~EUR avec $\beta=0.02$, 1\,840\,368~EUR avec $\beta=0.40$~: un écart de
511\,456~EUR.
\end{solution}

\begin{solution}{exo:rc:forward-rate-and-market-models:8}
Chaque forward n'est une martingale que sous sa propre mesure de paiement~; sous une
mesure de simulation commune, tous sauf un ont un drift
(voir la \cref{prop:rc:forward-rate-and-market-models:spotdrift}). Avec des drifts nuls, la
simulation fausse la valorisation de chaque caplet, sauf celui dont la mesure est
celle du numéraire, d'une erreur croissant avec la maturité et la volatilité, et
l'actualisation au compte spot est incohérente avec les forwards.
\end{solution}

\begin{solution}{pb:rc:forward-rate-and-market-models:1}
\textbf{1.} 2.933~\% et 4.102.
\textbf{2.} La variance intégrée de chaque forward jusqu'à sa fixation (le $\phi_k$
compte tenu de la forme de la volatilité)~; ni la corrélation, ni la façon dont la variance de chaque forward
se répartit dans le temps au-delà de ce qu'impose la forme retenue.
\textbf{3.} 33.0~\% et 23.1~\%.
\textbf{4.} Il dépend d'un seul forward, $F_5$.
\textbf{5.} Les paramètres de forme $a,b,c,d$, qui déplacent la variance entre les temps
précoces et tardifs, et modifient ainsi les swaptions d'échéances différentes et les volatilités forward.
\textbf{6.} 22.1~\% et 17.3~\%.
\textbf{7.} 2.35~millions d'euros et 1.84~millions d'euros.
\textbf{8.} 22~\%.
\textbf{9.} À environ deux erreurs types près (moins de 0.9 point de volatilité sur chaque
swaption d'échéance cinq ans).
\textbf{10.} Quelque chose comme l'estimation sur les Treasuries du chapitre~6, soit 0.77 entre les variations à deux et
à dix ans~: les taux forward distants de quelques années sont fortement mais non parfaitement
corrélés, ce qui suggère un $\beta$ intermédiaire entre les deux extrêmes.
\textbf{11.} Des swaptions sur des swaps plus longs (la \omterm{def:rc:sabr-in-rates-and-the-volatility-cube:matrix}{matrice des swaptions}), \omterm{def:rc:convexity-adjustments-and-constant-maturity-products:spread}{options sur spread CMS}
là où elles se négocient.
\textbf{12.} Environ 0.16~: avec $\beta=0.16$, la \omterm{def:rc:forward-rate-and-market-models:rebonato}{formule de Rebonato} donne 20~\%.
\textbf{13.} Toute la forme de la matrice de corrélation au-delà d'un seul paramètre, la
dépendance temporelle de la corrélation, et tout ce qui dépend de mouvements conjoints
à d'autres échéances.
\textbf{14.} Elle modifie la répartition des vegas des swaptions entre les forwards et les
rapports de couverture entre caps et swaptions~: un portefeuille couvert par des caps contre des swaptions
porte un risque de corrélation.
\textbf{15.} Un seul facteur fait bouger tous les taux ensemble~: corrélation parfaite, la limite
$\beta\to0$.
\textbf{16.} Le prix à un $\beta$ calibré sur la \omterm{def:rc:sabr-in-rates-and-the-volatility-cube:matrix}{matrice des swaptions} cotée, avec une
réserve couvrant une plage plausible de corrélation (ici l'écart entre
les valeurs de $\beta$ encadrant les niveaux historiques et implicites).
\textbf{17.} Suivre quotidiennement les corrélations calibrées et historiques, signaler les sauts,
et présenter séparément le P\&L issu de la réévaluation de $\beta$.
\textbf{18.} Comment il est choisi, sur quels instruments il est calibré, sa plage, les
produits qui y sont sensibles et la réserve constituée.
\textbf{19.} Résultat à retenir~: \emph{la corrélation qu'on ne voit pas}~: avec les mêmes caplets,
la swaption payeuse cinq ans sur cinq ans vaut 2.35~millions d'euros avec $\beta=0.02$ et 1.84
millions d'euros avec $\beta=0.40$, soit un écart de 22~\% (511\,456~EUR).
\textbf{20.} Comment des taux d'échéances différentes évoluent ensemble.
\end{solution}

\begin{solution}{iq:rc:forward-rate-and-market-models:1}
$\alpha(t,T)=\sigma_f(t,T)\int_t^T\sigma_f(t,u)\,du$ sous la probabilité risque-neutre. Elle
provient de l'exigence que tout prix d'obligation actualisé soit une martingale~: le
drift du log-prix, $r-\int\alpha+\frac12(\int\sigma_f)^2$, doit égaler $r$ pour toute
échéance~; dériver par rapport à $T$.
\iqlookfor{absence d'arbitrage sur les obligations, et la dérivation en deux lignes.}
\end{solution}

\begin{solution}{iq:rc:forward-rate-and-market-models:2}
Choisir une mesure (spot ou terminale), écrire le drift de chaque forward à partir de la
covariance avec le changement de numéraire, et faire évoluer $\ln F_k$ avec des normales corrélées,
un drift prédicteur-correcteur et le numéraire roulé à chaque fixation. Le drift est
là parce que chaque forward n'est une martingale que sous sa propre mesure de paiement.
\iqlookfor{mesure commune, drift issu du changement de numéraire, passage au logarithme.}
\end{solution}

\begin{solution}{iq:rc:forward-rate-and-market-models:3}
Un taux de swap est une moyenne de forwards pondérée par des poids aléatoires~; si les forwards sont
lognormaux sous leurs mesures, un taux de swap n'est pas lognormal sous sa probabilité annuité,
et inversement. Un des deux ensembles n'est évalué par Black que
de façon approximative.
\iqlookfor{les sommes de lognormales ne sont pas lognormales.}
\end{solution}

\begin{solution}{iq:rc:forward-rate-and-market-models:4}
Les caplets fixent la variance de chaque forward, mais pas leur corrélation, qui pilote
les volatilités des swaptions~: leurs prix sont fixés par une hypothèse de corrélation
arbitraire, pouvant aller jusqu'à 22~\% d'écart ici. Calibrer aussi sur les swaptions.
\iqlookfor{la corrélation comme donnée manquante.}
\end{solution}

\begin{solution}{iq:rc:forward-rate-and-market-models:5}
Hull--White~: rapide (arbres), ajustement exact aux swaptions co-terminales avec $\sigma(t)$, un seul facteur donc
des exercices parfaitement corrélés, ce qui sous-estime la valeur de l'option de bascule. Le modèle de
marché~: corrélation souple, extensions de smile, calibré sur les caplets et
les swaptions, mais demande une simulation de Monte-Carlo par régression pour l'exercice anticipé et est plus lent
et plus bruité. Utiliser Hull--White en production, le modèle de marché comme modèle de référence (benchmark).
\iqlookfor{rapidité contre corrélation, et partage production/référence.}
\end{solution}

\begin{solution}{iq:rc:forward-rate-and-market-models:6}
Biais de discrétisation du drift~: un drift d'Euler évalué au début des pas longs
est faux parce qu'il dépend de forwards qui bougent à l'intérieur du pas. Utiliser
un prédicteur-correcteur, des pas plus fins près des fixations, ou une mesure qui place
le forward sans drift du numéraire là où cela compte~; vérifier les caplets face à Black à
chaque taille de pas.
\iqlookfor{discrétisation du drift et contrôle de convergence.}
\end{solution}
